% !PRINTED 440
\weberpage{440}
%<<CONT
in $\frac{1}{2}(y - \sqrt{-m}\, x)$ auf, und diese beiden Größen müssen also relativ prim sein. Folglich kann die Gleichung (12) nicht bestehen.

Die quadratische Gleichung, deren beide Wurzeln $(lk)$, $(l^{-1}k)$ sind, hat also mit (18) % [sic] printed (18); no antecedent equation (18) occurs in this section here
\gesperrt{nur eine Wurzel} gemein, und daraus ergibt sich der wichtige Satz:

III. \gesperrt{Nach Adjunktion von $\sqrt{D}$ ist $(lk)$ rational ausdrückbar durch $(k)$ in der Form:}
\weq{13}{(lk) = f_l(k, \sqrt{D}),}
\gesperrt{wo $f_l$ eine rationale Funktion ist, deren Form durch die Klasse $l$ allein bestimmt ist.}

Die Änderung des Zeichens von $\sqrt{D}$ hat den Erfolg, daß $(lk)$ in $(l^{-1}k)$ übergeht, also:
\[ f_l(k, \sqrt{D}) = f_{l^{-1}}(k, -\sqrt{D}). \]
Ist $l$ zweiseitig, so gilt nach I. dieselbe Formel, nur daß $\sqrt{D}$ in $f_l$ dann nicht vorkommt.

Ist $l'$ eine zweite Klasse von derselben Beschaffenheit wie $l$, so kann man die Formel (13) auf $l'k$ anwenden und erhält:
\[ (l\, l' k) = f_l(l' k, \sqrt{D}), \]
was mit Anwendung von
\[ (l' k) = f_{l'}(k, \sqrt{D}) \]
in eine Gleichung der Form:
\[ (l\, l' k) = f_{l l'}(k, \sqrt{D}) \]
übergeht. Da man auf der linken Seite $l$ mit $l'$ vertauschen kann, so folgt
\[ f_{l l'}(k, \sqrt{D}) = f_{l' l}(k, \sqrt{D}), \]
so daß $f_{l l'}$ nur von der zusammengesetzten Klasse $l l'$ abhängt. Ebenso läßt sich ableiten
\[ (l\, l'\, l'' k) = f_{l l' l''}(k, \sqrt{-m}), \] % Der Druck wechselt hier von \sqrt{D} zu \sqrt{-m}; wegen D=-m sind die Bezeichnungen gleichwertig.
und da nach \S\,118, II. jede beliebige Klasse $s$ der Diskriminante $D$ aus solchen Klassen $l$ zusammensetzbar ist, so ist die Formel (13) nebst den daraus gezogenen Folgerungen nicht mehr an die über $l$ gemachte besondere Voraussetzung gebunden, daß der erste Koeffizient eine Primzahl sei.

\gesperrt{Wir haben also, wenn $s, k$ zwei beliebige Klassen der Diskriminante $D$ bedeuten, die Formeln:}
\weq{14}{(sk) = f_s(k, \sqrt{D}),}
\weq{15}{(s^{-1}k) = f_s(k, -\sqrt{D}).}
